\cotefile{58}
\scanpage{36}
\uline{Remarque 3} La condition $LG$ \unsure{et} $LG$ effective est stable par passage de $X$ à un voisinage ouvert $U$ de $Y$. Il n'en est plus de $=$ en général des conditions $(LG)_n$, \deleted{\unsure{en particulier|du point de vue} $(LF)_0$ Cependant on peut prouver facilement} qu'il en est \deleted{\unsure{encore} ainsi de $(LF)_0$, [et d'ailleurs} probablement \unsure{pas|que} $=$ pour $(LG)_0$, en vertu du fait que un module loc. libre $F$ sur un voisinage $U$ de $Y$ n'est pas restriction d'un module loc libre sur \deleted{un voisinage} $X$. Si cependant tt tel $F$ est \unsure{isom} à un \uline{quotient} d'un module loc. libre sur $X$ [p. ex. si tt module coh. sur $X$ est quotient d'un module loc libre, p. ex. $X$ quasi-projectif] \uline{alors} $(LG)_0$ pour $(X,Y)$ implique $(LG)_n$ pour tt $(U,Y)$ i.e. $LG$ pour $(X,Y)$. En effet, on est ramené à prouver que pour $F$ loc libre \uline{sur $X$}, $H^0(\unsure{U|X},F)\longrightarrow H^0(U_{/Y},F_{/Y})$ est bijectif, on sait déjà qu'il est surjectif, [car $H^0(X,F)\longrightarrow H^0(U,F)\longrightarrow H^0(U_{/Y},F_{/Y})$ l'est) reste à vérifier l'\uline{injectivité}, \deleted{\unsure{pour}}, \deleted{\unsure{du fait|donc}} \inserted{ce qui résulte de}
% archive sheet number: 35
